% ============================================================================
% CHAPTER 1: RESEARCH FUNDAMENTALS
% ============================================================================

\chapter{کلیات تحقیق}

\section{بيان مسأله}
در عصر دیجیتال، تبادل اطلاعات به‌ویژه در قالب تصاویر، به بخشی جدایی‌ناپذیر از ارتباطات، تجارت، سیستم‌های نظارتی، و رسانه‌های اجتماعی تبدیل شده است. تصاویر دیجیتال به دلیل ویژگی‌هایی همچون حجم زیاد، ساختار دوبعدی، همبستگی آماری بالا میان پیکسل‌های مجاور، و حساسیت بالا به تغییرات، به الگوریتم‌های رمزنگاری تخصصی‌تر و پیچیده‌تری نیاز دارند [1-3].

مهم‌ترین چالش‌ها در رمزنگاری تصاویر دیجیتال را می‌توان به شرح زیر برشمرد:
\begin{itemize}
    \item حفظ ویژگی‌های آماری تصویر پس از رمزنگاری؛
    \item ایجاد مقاومت در برابر انواع حملات نظیر تحلیل آماری، حملات تمایزی و جستجوی فراگیر (\lr{Brute Force})؛
    \item کاهش بار محاسباتی جهت کاربرد در سامانه‌های بلادرنگ؛
    \item پیشگیری از تخریب داده‌های تصویری طی فرآیند رمزنگاری و رمزگشایی.
\end{itemize}

در پاسخ به این چالش‌ها، بهره‌گیری از سیستم‌های آشوبی (\lr{Chaotic Systems}) به‌راهکاری مؤثر تبدیل شده است. با این حال، بروز پدیده‌ای به‌نام «تخریب آشوب» (\lr{Chaotic Degradation}) در سیستم‌های کلاسیک نظیر \lr{Logistic} و \lr{Tent}، امنیت سیستم را تهدید می‌کند [5]. سیستم‌های آشوبی نمایی (\lr{Exponential Chaotic Maps}) به‌عنوان نسل جدید، قادرند رفتارهای پیچیده‌تر و توزیع آماری یکنواخت‌تری ایجاد کنند [4].

\section{اهمیت و ضرورت تحقیق}
تصاویر رنگی به‌دلیل حجم بالای اطلاعات و کاربرد گسترده در حوزه‌های پزشکی و نظامی، بیش از گذشته در معرض تهدیدات قرار دارند. تصاویر رنگی به‌صورت ترکیبی از سه کانال \lr{RGB} تعریف می‌شوند که وابستگی بین کانال‌ها، کار رمزنگاری را دشوار می‌سازد [3]. طراحی الگوریتمی که بتواند این وابستگی‌ها را کاهش دهد، امری حیاتی است.

\section{اهداف تحقيق}
\subsection{هدف کلی}
طراحی و توسعه‌ی الگوریتمی برای رمزنگاری تصاویر رنگی با بهره‌گیری از ترکیب ۹ سیستم آشوبی نمایی و سیستم \lr{Chen}، با هدف ارتقاء امنیت و بهبود کارایی محاسباتی.

\subsection{اهداف جزئی}
\begin{itemize}
    \item بررسی عملکرد ۹ سیستم آشوبی نمایی در تولید دنباله‌های تصادفی پیچیده.
    \item بهبود مقاومت الگوریتم در برابر حملات آماری و جستجوی فراگیر.
    \item مقایسه امنیت و کارایی الگوریتم پیشنهادی با سایر روش‌های رمزنگاری.
\end{itemize}

\section{مفاهیم کلیدی تحقیق}
\textbf{رمزنگاری (\lr{Encryption}):} فرایندی است که طی آن اطلاعات به‌کمک کلید به صورتی غیرقابل‌فهم تبدیل می‌شود.

\textbf{سیستم آشوبی (\lr{Chaotic System}):} سیستم‌های دینامیکی غیرخطی که رفتار آن‌ها به شدت به شرایط اولیه وابسته است.

\textbf{تخریب آشوبی (\lr{Chaotic Degradation}):} از دست رفتن ویژگی‌های آشوبی در اثر خطای عددی یا پارامترهای خاص.

\section{سهم علمی پژوهش}
\begin{enumerate}
    \item ارائه الگوریتم رمزنگاری مبتنی بر ترکیب ۹ نقشه آشوبی نمایی متفاوت.
    \item طراحی سازوکار انتخاب دینامیک نقشه برای هر پیکسل.
    \item بهره‌گیری از سیستم آشوبی پیوسته \lr{Chen} برای جایگشت پیکسل‌ها.
    \item اثبات مقاومت الگوریتم در برابر حملات دیفرانسیلی و آماری.
\end{enumerate}

\section{ساختار پایان‌نامه}
این پایان‌نامه در پنج فصل سازمان‌دهی شده است. فصل دوم مفاهیم پایه و پیشینه را بررسی می‌کند. فصل سوم الگوریتم پیشنهادی را معرفی می‌کند. فصل چهارم پیاده‌سازی و نتایج تجربی را ارائه می‌دهد. و فصل پنجم نتایج را تحلیل و مقایسه می‌کند.
\
